\section{2019年高考题}
\subsection{2019年全国I卷（理）}\label{gk:2019年全国I卷（理）}


\clearpage


\subsection{2019年全国I卷（文）}\label{gk:2019年全国I卷（文）}


\clearpage


\subsection{2019年全国II卷（理）}\label{gk:2019年全国II卷（理）}


\clearpage


\subsection{2019年全国II卷（文）}\label{gk:2019年全国II卷（文）}


\clearpage


\subsection{2019年全国III卷（理）}\label{gk:2019年全国III卷（理）}


\clearpage


\subsection{2019年全国III卷（文）}\label{gk:2019年全国III卷（文）}


\clearpage


\subsection{2019年北京卷（理）}\label{gk:2019年北京卷（理）}

\begin{description}
    \item[一、]
    \begin{enumerate}
        \item 1
        \item 2
        \item 3
        \item 4
        \item 5
        \item 6
        \item 7
        \item 8
    \end{enumerate}
    \item[二、]
    \begin{enumerate}
      \addtocounter{enumi}{+8}   
        \item 1
        \item 2
        \item 3
        \item 4
        \item 1
        \item 6
    \end{enumerate}
    \item[三、]   
    \begin{enumerate}
      \addtocounter{enumi}{+14}   
        \item 1
        \item 2
        \item 3
        \item 已知抛物线 $C: x^{2} = -2py$ 经过点 $(2,-1)$。

\begin{enumerate}[label=(\arabic*)]
    \item 求抛物线 $C$ 的方程及其准线方程；
    \item 设 $O$ 为原点，过抛物线 $C$ 的焦点作斜率不为 $0$ 的直线 $l$ 交抛物线 $C$ 于两点 $M$, $N$，直线 $y = -1$ 分别交直线 $OM$, $ON$ 于点 $A$ 和点 $B$。求证：以 $AB$ 为直径的圆经过 $y$ 轴上的两个定点。
\end{enumerate}

        \item 已知函数 $f(x) = \frac{1}{4}x^{3} - x^{2} + x$。

\begin{enumerate}[label=(\arabic*)]
    \item 求曲线 $y = f(x)$ 的斜率为 $1$ 的切线方程；
    \item 当 $x \in [-2,4]$ 时，求证：$x - 6 \leq f(x) \leq x$；
    \item 设 $F(x) = |f(x) - (x + a)|$ ($a \in \mathbb{R}$)，记 $F(x)$ 在区间 $[-2,4]$ 上的最大值为 $M(a)$。当 $M(a)$ 最小时，求 $a$ 的值。
\end{enumerate}
        \item 已知数列 $\{a_n\}$，从中选取第 $i_1$ 项、第 $i_2$ 项、$\cdots$、第 $i_m$ 项 ($i_1 < i_2 < \cdots < i_m$)，若 $a_{i_1} < a_{i_2} < \cdots < a_{i_m}$，则称新数列 $a_{i_1}, a_{i_2}, \cdots, a_{i_m}$ 为 $\{a_n\}$ 的长度为 $m$ 的递增子列。规定：数列 $\{a_n\}$ 的任意一项都是 $\{a_n\}$ 的长度为 $1$ 的递增子列。

\begin{enumerate}[label=(\arabic*)]
    \item 写出数列 $1, 8, 3, 7, 5, 6, 9$ 的一个长度为 $4$ 的递增子列；
    \item 已知数列 $\{a_n\}$ 的长度为 $p$ 的递增子列的末项的最小值为 $a_{m_0}$，长度为 $q$ 的递增子列的末项的最小值为 $a_{n_0}$。若 $p < q$，求证：$a_{m_0} < a_{n_0}$；
    \item 设无穷数列 $\{a_n\}$ 的各项均为正整数，且任意两项均不相等。若 $\{a_n\}$ 的长度为 $s$ 的递增子列末项的最小值为 $2s - 1$，且长度为 $s$ 末项为 $2s - 1$ 的递增子列恰有 $2^{s-1}$ 个 ($s = 1, 2, \cdots$)，求数列 $\{a_n\}$ 的通项公式。
\end{enumerate}
    \end{enumerate}
\end{description}
\clearpage


\subsection{2019年北京卷（文）}\label{gk:2019年北京卷（文）}

\begin{description}
    \item[一、]
    \begin{enumerate}
        \item 1
        \item 2
        \item 3
        \item 4
        \item 5
        \item 6
        \item 7
        \item 8
    \end{enumerate}
    \item[二、]
    \begin{enumerate}
      \addtocounter{enumi}{+8}   
        \item 1
        \item 2
        \item 3
        \item 4
        \item 1
        \item 6
    \end{enumerate}
    \item[三、]   
    \begin{enumerate}
      \addtocounter{enumi}{+14}   
        \item 1
        \item 2
        \item 3
        \item 
        \item 已知椭圆 $C: \frac{x^{2}}{a^{2}} + \frac{y^{2}}{b^{2}} = 1$ 的右焦点为 $(1,0)$，且经过点 $A(0,1)$。

\begin{enumerate}[label=(\arabic*)]
    \item 求椭圆 $C$ 的方程；
    \item 设 $O$ 为原点，直线 $l: y = kx + t$ ($t \neq \pm 1$) 与椭圆 $C$ 交于两个不同点 $P$, $Q$，直线 $AP$ 与 $x$ 轴交于点 $M$，直线 $AQ$ 与 $x$ 轴交于点 $N$，若 $|OM| \cdot |ON| = 2$，求证：直线 $l$ 经过定点。
\end{enumerate}
        \item 已知函数 $f(x) = \frac{1}{4}x^{3} - x^{2} + x$。

\begin{enumerate}[label=(\arabic*)]
    \item 求曲线 $y = f(x)$ 的斜率为 $1$ 的切线方程；
    \item 当 $x \in [-2,4]$ 时，求证：$x - 6 \leq f(x) \leq x$；
    \item 设 $F(x) = |f(x) - (x + a)|$ ($a \in \mathbb{R}$)，记 $F(x)$ 在区间 $[-2,4]$ 上的最大值为 $M(a)$。当 $M(a)$ 最小时，求 $a$ 的值。
\end{enumerate}
    \end{enumerate}
\end{description}
\clearpage


\subsection{2019年江苏卷}\label{gk:2019年江苏卷}


\clearpage


\subsection{2019年上海卷}\label{gk:2019年上海卷}


\clearpage


\subsection{2019年天津卷（理）}\label{gk:2019年天津卷（理）}


\clearpage


\subsection{2019年天津卷（文）}\label{gk:2019年天津卷（文）}


\clearpage


\subsection{2019年浙江卷}\label{gk:2019年浙江卷}




\clearpage